\documentclass{beamer}
\usetheme{Madrid}
\usepackage[utf8]{inputenc}
\usepackage[T1]{fontenc}
\usepackage[french]{babel}
\usepackage{amsmath}
\usepackage{tikz}

\title{Le Perceptron Multicouche (MLP)}
\subtitle{Une courte introduction}
\author{}
\date{\today}

\begin{document}

\begin{frame}
  \titlepage
\end{frame}

\begin{frame}{Qu'est-ce qu'un MLP ?}
  \begin{itemize}
    \item Un \textbf{réseau de neurones à propagation avant} (\emph{feedforward}).
    \item Composé de plusieurs couches de neurones :
    \begin{itemize}
      \item une couche d'\textbf{entrée},
      \item une ou plusieurs couches \textbf{cachées},
      \item une couche de \textbf{sortie}.
    \end{itemize}
    \item Chaque neurone d'une couche est connecté à tous les neurones de la couche suivante (couches \emph{denses}).
    \item Utilisé en classification, en régression, et comme brique de base de nombreux modèles modernes.
  \end{itemize}
\end{frame}

\begin{frame}{Architecture}
  \centering
  \begin{tikzpicture}[
      neuron/.style={circle, draw, minimum size=0.7cm},
      >=stealth, scale=1.1]
    % Entrée
    \foreach \i in {1,2,3}
      \node[neuron, fill=blue!20] (I\i) at (0,-\i) {$x_\i$};
    % Cachée
    \foreach \j in {1,2,3,4}
      \node[neuron, fill=green!20] (H\j) at (2.5,-\j+0.5) {};
    % Sortie
    \foreach \k in {1,2}
      \node[neuron, fill=red!20] (O\k) at (5,-\k-0.5) {$\hat{y}_\k$};
    % Connexions
    \foreach \i in {1,2,3}
      \foreach \j in {1,2,3,4}
        \draw[->, gray] (I\i) -- (H\j);
    \foreach \j in {1,2,3,4}
      \foreach \k in {1,2}
        \draw[->, gray] (H\j) -- (O\k);
    % Légendes
    \node at (0,0) {\small Entrée};
    \node at (2.5,0.9) {\small Couche cachée};
    \node at (5,-0.6) {\small Sortie};
  \end{tikzpicture}
\end{frame}

\begin{frame}{Fonctionnement d'un neurone}
  Chaque couche $l$ calcule :
  \[
    \mathbf{h}^{(l)} = \sigma\!\left( W^{(l)} \mathbf{h}^{(l-1)} + \mathbf{b}^{(l)} \right)
  \]
  \begin{itemize}
    \item $W^{(l)}$ : matrice des poids, $\mathbf{b}^{(l)}$ : biais.
    \item $\sigma$ : fonction d'activation \textbf{non linéaire} :
    \begin{itemize}
      \item ReLU : $\sigma(z) = \max(0, z)$
      \item Sigmoïde : $\sigma(z) = \dfrac{1}{1 + e^{-z}}$
      \item Tangente hyperbolique : $\tanh(z)$
    \end{itemize}
    \item Sans non-linéarité, le MLP se réduirait à un simple modèle linéaire.
  \end{itemize}
\end{frame}

\begin{frame}{Apprentissage}
  \begin{enumerate}
    \item \textbf{Propagation avant} : calcul de la prédiction $\hat{y}$.
    \item \textbf{Fonction de perte} : mesure l'erreur, par ex.
      \[
        \mathcal{L} = \frac{1}{N}\sum_{i=1}^{N} (y_i - \hat{y}_i)^2
        \quad \text{ou l'entropie croisée.}
      \]
    \item \textbf{Rétropropagation} : calcul des gradients $\frac{\partial \mathcal{L}}{\partial W}$ par la règle de dérivation en chaîne.
    \item \textbf{Descente de gradient} : mise à jour des poids
      \[
        W \leftarrow W - \eta \, \frac{\partial \mathcal{L}}{\partial W}
      \]
      où $\eta$ est le taux d'apprentissage.
  \end{enumerate}
\end{frame}

\begin{frame}{Forces et limites}
  \begin{columns}[T]
    \begin{column}{0.48\textwidth}
      \textbf{Forces}
      \begin{itemize}
        \item Approximateur universel de fonctions.
        \item Simple à implémenter.
        \item Polyvalent.
      \end{itemize}
    \end{column}
    \begin{column}{0.48\textwidth}
      \textbf{Limites}
      \begin{itemize}
        \item Risque de surapprentissage.
        \item Nombreux paramètres.
        \item N'exploite pas la structure des données (images, séquences) : on lui préfère alors les CNN ou les Transformers.
      \end{itemize}
    \end{column}
  \end{columns}
\end{frame}

\begin{frame}{Conclusion}
  \begin{itemize}
    \item Le MLP est le modèle fondamental de l'apprentissage profond.
    \item Il combine couches denses et activations non linéaires.
    \item Il s'entraîne par rétropropagation et descente de gradient.
    \item Il reste une composante clé des architectures modernes.
  \end{itemize}
  \vfill
  \centering
  \Large Merci de votre attention !
\end{frame}

\end{document}
